Here we discuss works that are closely related to the framework and sampler modifications proposed in this work.

\subsection{Adaptive Block Samplers}
Our work is not the first to suggest incrementally moving the denoising window: \citet{luo2026dsb} applies this to masked diffusion and sees performance gains.
\citet{gat2025set} fine-tunes an auto-regressive model to simultaneously decode the masked tokens ahead using the EB sampler of \citet{ben2026accelerated} while generating the next token auto-regressively, which is the same core idea of a rolling denoising window.
Furthermore, the denoising window is entirely defined by the entropy bound, which has been already defined in \citet{ben2026accelerated} and used in the original sampler of \citet{team2026diffusiongemma}.

Similar ideas of combining diffusion with the left-to-right bias of language have appeared even earlier, such as in \citep{wu2023ar}.
Our contribution is a study of the tradeoff between per-position commitment and block-level commitment, using DiffusionGemma as the experimental setting. Frozen EB simulates the loss of later revision under per-position commitment. RES waits for a contiguous accepted prefix, preserving sampler flexibility for accepted positions farther ahead while avoiding redundant computation on the prefix. The monotone masked-diffusion rules in Table~\ref{tab:samplers} provide related examples of preventing later revision on a per-position basis, though freezing does not necessarily imply immediate KV caching.

\subsection{Rolling at Training Time: Clock Diffusion}
Concurrent with our work, \citet{schiff2026clock} introduce Clock Diffusion, which develops essentially the same idea for Gaussian diffusion: a window that rolls left to right and commits finished positions to the KV cache, with a position-dependent noise schedule that mirrors our position-based temperature schedule.

\subsection{Rolling Diffusion Beyond Text}
Beyond text generation, the idea of rolling diffusion has been studied extensively in other domains.
Within video generation, rolling diffusion was a natural mechanism to address the temporal nature of video generation and avoid performing diffusion on every frame simultaneously \citep{ruhe2024rolling}.
Many other works include similar mechanisms \citep{chen2024diffusion, song2025history, xie2025progressive}, and \citet{papa2026asyncpatch} carry the idea over to images by assigning different patches different levels of corruption.
